\section{Conclusion}
In this work, we introduced MCPTox, the first benchmark designed to systematically evaluate Tool Poisoning Attacks within the real-world MCP ecosystem. We evaluated 20 prominent LLM agents and conducted a comprehensive analysis, revealing a widespread vulnerability where many popular LLM agents exhibited attack success rates exceeding 60\%, with the highest ASR reaching 72\%. This work establishes Tool Poisoning as a practical threat, underscores the risks associated with how agents trust tool metadata, and emphasizes an urgent need for a pre-execution security mechanism to safeguard against these attacks.
\section{Ethical Considerations}
The primary ethical consideration of our research in developing the MCPTox benchmark stems from the dual-use nature of disclosed vulnerabilities. By revealing these vulnerabilities, our goal is to preemptively strengthen the community against potential exploits, thereby promoting a culture of enhanced security and resilience. Although we acknowledge that sharing information about these weaknesses might lead to their misuse, we argue that it is crucial to be aware of them in order to safeguard against such threats. 
% Therefore, we believe that our paper is in alignment with ethical principles. Additionally, we have disclosed our findings to the developers of the evaluated models to ensure they are aware of the vulnerabilities. 